\section{Light-Cone Coordinates}
\label{ap:Light_Cone}
\textcolor{red}{May need to be reviewed so it has a similar structure and flow to the Bjorken coordinates appendix.}

Throughout the derivations, it will be useful to work in light-cone coordinates. A four-vector $a^\mu = (a^0,a^1,a^2,a^3)$ may be expressed in light-cone coordinates as
\begin{equation}
a^\mu = (a^+, a^-, a^2, a^3) \equiv (a^+, a^-, \mathbf{a}_\perp), \qquad \mu = (+,-,~2,~3)
\end{equation}

\noindent with
\begin{equation}
	a^+ = \frac{a^0 + a^1}{\sqrt{2}} \qquad a^- = \frac{a^0 - a^1}{\sqrt{2}} \qquad \mathbf{a}_\perp = (a^2,a^3)
\end{equation}

The dot product between two four-vectors $a^\mu$ and $b^\nu$ becomes
\begin{equation}
	a \cdot b =  a^+ b^- + a^- b^+ - \mathbf{a}_\perp \cdot \mathbf{b}_\perp = a^+ b^- + a^- b^+ - a^2 b^2 - a^3 b^3
\end{equation}

\noindent the metric being
\begin{equation}
	g_ {\mu \nu} = \begin{bmatrix}
		0 & 1 & 0 & 0 \\
		1 & 0 & 0 & 0 \\
		0 & 0 & -1 & 0 \\
		0 & 0 & 0 & -1
	\end{bmatrix}
	\label{eq:light_cone_metric}
\end{equation}

The components of the corresponding covariant four-vector $a_\mu$ are
\begin{equation}
	a_\mu = (a_+, a_-, a_2, a_3) = (a^-, a^+, -a^2, -a^3)
\end{equation}

Considering the position four-vector,
\begin{equation}
	x^\mu = (x^+, x^-, x^2, x^3) = \left( \frac{x^0+x^1}{\sqrt{2}}, \frac{x^0-x^1}{\sqrt{2}}, x^2, x^3  \right)
\end{equation}
we may obtain the new base vectors. Focusing on the first two coordinates:
\begin{equation}
	\vec{r} = x^0 \vec{e}_0 + x^1 \vec{e}_1, \quad \text{with} \quad \vec{e}_0 = (1,0,0,0), \quad \vec{e}_1=(0,1,0,0)
\end{equation}
\begin{equation}
\begin{split}
	\vec{e}_+ & = \frac{\partial \vec{r}}{\partial x^+} = \frac{1}{\sqrt{2}} \vec{e}_0 + \frac{1}{\sqrt{2}} \vec{e}_1 \\
	\vec{e}_- & = \frac{\partial \vec{r}}{\partial x^-} = \frac{1}{\sqrt{2}} \vec{e}_0 - \frac{1}{\sqrt{2}} \vec{e}_1
\end{split}
\end{equation}

Allowing us to rewrite $\vec{r}$:
\begin{equation}
	\vec{r} =  x^+ \vec{e}_+ + x^- \vec{e}_-
\end{equation}

Thus, the $x^+$ and $x^-$ axes lie on the light cone (Fig. \ref{fig:light_cone}).


\begin{figure}[h]\centering
  \includetikz{Diagrams/light_cone_diagram}
	\caption{Position four-vector axes in light-cone coordinates.}
	\label{fig:light_cone}
\end{figure}

As for the four-gradient,
\begin{equation}
	\partial_\mu = (\partial_+, \partial_-, \partial_2, \partial_3) = \left( \frac{\partial}{\partial x^+}, \frac{\partial}{\partial x^-}, \frac{\partial}{\partial x^2}, \frac{\partial}{\partial x^3} \right)
\end{equation}

\begin{equation}
	\partial^\mu = (\partial^+, \partial^-, \partial^2, \partial^3) = (\partial_-, \partial_+, -\partial_2, -\partial_3)
\end{equation}

%\textcolor{red}{Conveniently, the Dirac matrices (in the Dirac representation?) "in light cone coordinates" have anticommutation relations obeying to our new metric.}

\textcolor{red}{Abusing language and the notation,} one may similarly define the Dirac matrices ``in light-cone coordinates".
\begin{equation}
	\gamma^+ = \frac{\gamma^0 + \gamma^1}{\sqrt{2}} \qquad \gamma^- = \frac{\gamma^0 - \gamma^1}{\sqrt{2}}
\end{equation}

They obey anticommutation relations of the same form as the usual Dirac matrices, with the new metric of Eq. \ref{eq:light_cone_metric}.

\begin{equation}
	\{ \gamma^\mu, \gamma^\nu \} = 2g^{\mu \nu}, \qquad \mu,\nu = (+,-,2,3)
\end{equation}

%\tdplotsetmaincoords{45}{125}
%\begin{tikzpicture}[tdplot_main_coords]
%\draw[thick,->] (-2,0,0) -- (2,0,0) node[anchor=north east]{$x$};
%\draw[thick,->] (0,-2,0) -- (0,2,0) node[anchor=north west]{$y$};
%\draw[thick,->] (0,0,-2) -- (0,0,2) node[anchor=south]{$z$};
%\end{tikzpicture}

\section{Bjorken/Milne Coordinates}
\label{ap:bjorken_coordinates}
\footnote{\textcolor{red}{what Margaret calls them on page 12}}
In this system, two familiar quantities are used as the first two coordinates: proper time, $\tau$ \textcolor{red}{of an object with $\mathbf{x}_\perp=0$? which set off from the origin with constant velocity?} -- which is invariant under boosts along the longitudinal direction, ($x^1$) --, and rapidity \textcolor{red}{of some body in some particular circumstances too}, $\eta$. The transverse coordinates are unchanged, as in light-cone coordinates.
\begin{equation}
\begin{split}
	\tau & = \sqrt{(x^0)^2 - (x^1)^2} = \sqrt{2x^+ x^-} \\
	\eta & = \frac{1}{2} \log \left( \frac{x^0+x^1}{x^0-x^1} \right) = \frac{1}{2} \log \left(\frac{x^+}{x^-} \right)
\end{split}
\end{equation}

As for the inverse transformation,
\begin{equation}
	x^+  = \frac{\tau}{\sqrt{2}} e^\eta \qquad
	x^-  = \frac{\tau}{\sqrt{2}} e^{-\eta}
\end{equation}

\begin{figure}[h]\centering
  \includetikz{Diagrams/bjorken_coordinates_diagram}

  \caption{Contour lines for proper time, $\tau$, in red, and rapidity, $\eta$, in blue, in the first quadrant. }
  \label{fig:bjorken_coordinates}
\end{figure}

We may find the new basis vectors: \textcolor{red}{versores n deviam ser "e"}
\begin{equation*}
	\vec{r} =  x^+ \vec{e}_+ + x^- \vec{e}_-
\end{equation*}
\begin{equation}
\begin{split}
	\vec{e}_\tau & = \frac{\partial \vec{r}}{\partial \tau} = \frac{e^\eta}{\sqrt{2}} \vec{e}_+ + \frac{e^{-\eta}}{\sqrt{2}} \vec{e}_- = \frac{x^+}{\tau} \vec{e}_+ + \frac{x^-}{\tau} \vec{e}_- \\
	\vec{e}_\eta & = \frac{\partial \vec{r}}{\partial \eta} = \frac{\tau}{\sqrt{2}} e^\eta \vec{e}_+ - \frac{\tau}{\sqrt{2}}e^{-\eta} \vec{e}_- = x^+ \vec{e}_+ - x^- \vec{e}_-
\end{split}
\end{equation}
\noindent which are not, in general, of unit length.

By requiring
\begin{equation*}
	\vec{A} = A^+ \vec{e}_+ + A^- \vec{e}_- \equiv A^\tau \vec{e}_\tau + A^\eta \vec{e}_\eta,
\end{equation*}
\noindent then
\begin{equation}
\begin{split}
	A^\tau & = \frac{x^- A^+ + x^+ A^-}{\tau} \\
	A^\eta & = \frac{x^- A^+ - x^+ A^-}{\tau^2}
\end{split}
\end{equation}

Particularly,
\begin{equation}
	\vec{r} = \tau \vec{e}_\tau
\end{equation}

\color{purple}
Additionally, under a boost of rapidity $\sigma$, the coordinates $\tau$ and $\eta$ are changed as:
\begin{subequations}
\begin{gather}
	\tau' = \tau \\
	\eta' = \eta - \sigma
\end{gather}
\end{subequations}
\color{black}
